\subsubsection*{4.19.}
从特殊情形 4.18 推出 4.15 和 4.17。设 $u:Y\longrightarrow X$ 是 $(X,X_0)$ 的一个消解，即诱导同构 $Y^*\xrightarrow{\ \sim\ }X^*$ 的固有态射，使 $Y$ 光滑且 $(Y_0)_{\mathrm{red}}$ 为 $Y$ 中的正规交叉除子。若 $V$ 沿 $X_0$ 拟幺幂，则其逆像 $u^*V$ 沿 $Y_0$ 拟幺幂。

对 $X$ 上的每个凝聚层 $F$，用 $F^m$ 或 $j_*^m j^*F$ 表示 $j_*j^*F$ 中由亚纯截面构成的子层。若 $f$ 为 $X_0$ 的一个方程，则 $j_*^m j^*F$ 是系统
\begin{equation}\tag{4.19.1}
F\xrightarrow{f}F\xrightarrow{f}F\longrightarrow\cdots.
\end{equation}
的归纳极限。以下引理的证明见 [2]。

\subsubsection*{引理 4.19.2.}
\underline{对 $X$ 上的每个凝聚层 $F$，有}
\[
F^m\xrightarrow{\ \sim\ }u_*(u^*F)^m
\]
\underline{且对 $i>0$ 有 $R^iu_*(u^*F)^m=0$。}

由 Grauert 定理，各 $R^iu_*u^*F$ 为凝聚层。$F\longrightarrow u_*u^*F$ 的核与余核，以及 $i>0$ 时的 $R^iu_*u^*F$，在 $X_0$ 外均为零，因而局部地被 $f$ 的某个幂零化。由于 $u$ 固有，$R^iu_*$ 与滤过归纳极限交换，特别地与归纳极限 (4.19.1) 交换，由此得到 (4.19.2)。

由 4.19.2 与 4.12 推出：在上述假设下，仍用 $u$ 表示 $u$ 诱导的映射 $Y_0\to X_0$，则有下述结论。

\subsubsection*{引理 4.19.3.}
\underline{态射}
\[
\Psi_\eta^m(F)\longrightarrow
u_*\Psi_\eta^m(u^*F)\longrightarrow
Ru_*\Psi_\eta^m(u^*F)
\]
\underline{都是同构。}

第二个箭头为同构意味着
\[
R^iu_*\Psi_\eta^m(u^*F)=0\qquad(i>0).
\]
这由 4.12 及 $i>0$ 时 $R^iu_*(i^*u^*F^m)=0$ 推出；后者又由 $i>0$ 时 $R^iu_*(u^*F^m)=0$ 和 $u$ 的固有基变换定理推出。

为证明 4.15，只需考虑交换图
\begin{center}
\begin{tikzcd}[
  column sep=large,
  row sep=large,
  cells={nodes={font=\small}}
]
\Psi^m(\Omega_{X^*}^*(V))
  \arrow[d]
  \arrow[r,equal,"(4.19.3)"']
& u_*\Psi^m(\Omega_{Y^*}^*(u^*V))
  \arrow[r,"\sim","(4.19.3)"']
& Ru_*\Psi^m(\Omega_{Y^*}^*(u^*V))
  \arrow[d,"\wr"',"(4.18)"]
\\
R\Psi(V)
  \arrow[rr,"(1.3.1;\ \mathrm{XIII}\ 2.1.7.1)"]
& {}
& Ru_*R\Psi(u^*V)
\end{tikzcd}
\end{center}
4.17 的证明相同，只需将所有 $\Psi^m$ 换为 $\Psi^{mqu}$。

\subsubsection*{注 4.20.}
定理 4.13 给出了单值定理在复形层次上的一个体现。

\subsection*{参考文献}

\begin{enumerate}[label={[\arabic*]},leftmargin=3em]
\item P. Deligne, \emph{Equations diff\'erentielles \`a points
  singuliers r\'eguliers}, Lecture Notes in Mathematics 163,
  Springer-Verlag.
\item A. Grothendieck, \emph{On the De Rham cohomology of algebraic
  varieties}, Publ. Math. IHES \underline{29}, 95--103 (1966).
\item K. Lamotke, \emph{Isolated critical points and monodromy}.
\end{enumerate}
